\chapter{Exemplo de Contagem e Exibição da Paginação}
\label{ap:paginacao}

O quadro mostra um trabalho hipotético. A quantidade real de páginas pré-textuais varia conforme os elementos utilizados.

\begin{longtable}{p{0.34\linewidth}p{0.20\linewidth}p{0.18\linewidth}p{0.16\linewidth}}
\toprule
Elemento & Participa da contagem? & Número lógico & Número visível? \\
\midrule
Capa & Não & -- & Não \\
Folha de rosto & Sim & 1 & Não \\
Dados catalográficos & Não & -- & Não \\
Folha de aprovação & Sim & 2 & Não \\
Dedicatória & Sim & 3 & Não \\
Agradecimentos & Sim & 4 & Não \\
Resumo & Sim & 5 & Não \\
Resumo estrangeiro & Sim & 6 & Não \\
Lista de ilustrações & Sim & 7 & Não \\
Sumário & Sim & 8 & Não \\
Primeira página da introdução & Sim & 9 & Sim: 9 \\
Demais textuais e pós-textuais & Sim & Sequencial & Sim \\
\bottomrule
\end{longtable}

\caixaatencao{Não se imprime numeração romana nos elementos pré-textuais do TCC. A tabela usa “número lógico” apenas para explicar a contagem que será exibida a partir da primeira página textual.}
